% ─────────────────────────────────────────────────────────────────────────────
% Environment: Containerlab (Virtual)
% 3 models × 10 runs + 1 reference (static scenario)
% Evaluation criteria: see §\ref{subsec:criteria}
% ─────────────────────────────────────────────────────────────────────────────
\section{Virtual Environment (Containerlab)}\label{sec:eval-containerlab}

\paragraph{Overview}

\begin{table}[H]
    \caption{Benchmark Results --- Containerlab, Single-Agent ($n = 10$ runs each)}
    \centering

    \resizebox{\textwidth}{!}{
        \begin{tabular}{l c c c c c}
            \hline\hline
            Model              & Avg.\ Duration & Avg.\ Tokens & Avg.\ Hosts & Avg.\ Services & Avg.\ Findings \\ [0.5ex]
            \hline
            claude-opus-4-7    & 122 s          & 52,888       & 5           & 9              & 8              \\
            qwen3:30b          & 362 s          & 86,446       & 4           & 7              & 5              \\
            gpt-oss:120b (BFH) & 963 s          & 62,475       & 5           & 6              & 4              \\
            Static Scenario    & 29 s           & ---          & 5           & 13             & 8              \\
            \hline
        \end{tabular}
    }
    \label{tab:overview-containerlab-single}
\end{table}

\begin{table}[H]
    \caption{Benchmark Results --- Containerlab, Multi-Agent ($n = 10$ runs each)}
    \centering

    \resizebox{\textwidth}{!}{
        \begin{tabular}{l c c c c c}
            \hline\hline
            Model              & Avg.\ Duration & Avg.\ Tokens & Avg.\ Hosts & Avg.\ Services & Avg.\ Findings \\ [0.5ex]
            \hline
            claude-opus-4-7    & 150 s          & 59,988       & 5           & 9              & 8              \\
            qwen3:30b          & 252 s          & 27,789       & 4           & 8              & 5              \\
            gpt-oss:120b (BFH) & 847 s          & 46,736       & 5           & 7              & 5              \\
            Static Scenario    & 29 s           & ---          & 5           & 13             & 8              \\
            \hline
        \end{tabular}
    }
    \label{tab:overview-containerlab-multi}
\end{table}

Tables~\ref{tab:overview-containerlab-single} and~\ref{tab:overview-containerlab-multi} summarise
the quantitative results across ten runs per model in the virtual \gls{Containerlab} environment,
for the single-agent and multi-agent configurations respectively.
The frontier model \texttt{claude-opus-4-7} achieves the shortest mean duration in single-agent mode
(122\,s) and only slightly longer in multi-agent mode (150\,s), while consistently attaining the
highest discovery quality, matching the Static Scenario baseline on findings (8 out of 8) and
recovering the most services (9) in both configurations.
The open-weight local model \texttt{qwen3:30b} benefits significantly from the multi-agent setup:
duration drops from 362\,s to 252\,s and token consumption falls from 86,446 to 27,789, while
service coverage improves from 7 to 8.
The BFH-hosted model \texttt{gpt-oss:120b} improves from 963\,s to 847\,s in multi-agent mode,
though it remains the slowest model across both configurations.
The \emph{Static Scenario} executes a fixed, deterministic reconnaissance script in 29\,s and serves
as the reference baseline: it discovers all 5 reachable hosts and finds 13 services (through exhaustive
port scanning) and 8 findings, setting the benchmark and reference for the reconnaissance scenarios
in this environment.

\subsection{Quantitative Criteria}\label{subsed:quantitative-criteria-containerlab}

\begin{figure}[H]
    \begin{subfigure}{0.5\linewidth}
        \centering
        \includesvg[width=1\linewidth]{figures/benchmark_duration_comparison_25_05_2026.svg}
        \caption{Single-agent.}
        \label{fig:benchmark_duration_containerlab_single}
    \end{subfigure}%
    \begin{subfigure}{0.5\linewidth}
        \centering
        \includesvg[width=1\linewidth]{figures/benchmark_duration_comparison_26_05_2026.svg}
        \caption{Multi-agent.}
        \label{fig:benchmark_duration_containerlab_multi}
    \end{subfigure}
    \caption{Mean run duration per model (Containerlab, $n=10$ runs each).}
    \label{fig:benchmark_duration_containerlab}
\end{figure}

Figure~\ref{fig:benchmark_duration_containerlab} visualises the mean run duration per model for
both configurations side by side.
In single-agent mode, \texttt{claude-opus-4-7} completes in 122\,s, while \texttt{qwen3:30b}
takes 362\,s and \texttt{gpt-oss:120b} 963\,s.
The multi-agent configuration reduces duration for all open-weight models:
\texttt{qwen3:30b} drops to 252\,s and \texttt{gpt-oss:120b} to 847\,s, while
\texttt{claude-opus-4-7} remains stable at 150\,s.
The Static Scenario at 29\,s reflects only the fixed execution time of the deterministic script
with no inference overhead.


\begin{figure}[H]
    \begin{subfigure}{0.5\linewidth}
        \centering
        \includesvg[width=1\linewidth]{figures/benchmark_token_usage_25_05_2026.svg}
        \caption{Single-agent.}
        \label{fig:benchmark_tokens_containerlab_single}
    \end{subfigure}%
    \begin{subfigure}{0.5\linewidth}
        \centering
        \includesvg[width=1\linewidth]{figures/benchmark_token_usage_26_05_2026.svg}
        \caption{Multi-agent.}
        \label{fig:benchmark_tokens_containerlab_multi}
    \end{subfigure}
    \caption{Mean total token usage per model (Containerlab, $n=10$ runs each).}
    \label{fig:benchmark_tokens_containerlab}
\end{figure}

Figure~\ref{fig:benchmark_tokens_containerlab} shows mean token consumption per model.
The most striking difference between configurations is \texttt{qwen3:30b}: in single-agent mode it
consumes the most tokens of any \acrshort{AI} model (86,446), whereas in multi-agent mode its
consumption falls to 27,789 --- the lowest of all three models.
This indicates that distributing the task across specialised sub-agents allows each agent to operate
on a much smaller context.
\texttt{claude-opus-4-7} uses 52,888 tokens in single-agent and 59,988 in multi-agent mode;
\texttt{gpt-oss:120b} drops from 62,475 to 46,736.
Token count alone is not a reliable proxy for result quality: \texttt{claude-opus-4-7} achieves the
best outcomes while consuming a moderate number of tokens in both configurations.


\begin{figure}[H]
    \begin{subfigure}{0.5\linewidth}
        \centering
        \includesvg[width=1\linewidth]{figures/benchmark_tokens_vs_duration_25_05_2026.svg}
        \caption{Single-agent.}
        \label{fig:benchmark_scatter_containerlab_single}
    \end{subfigure}%
    \begin{subfigure}{0.5\linewidth}
        \centering
        \includesvg[width=1\linewidth]{figures/benchmark_tokens_vs_duration_26_05_2026.svg}
        \caption{Multi-agent.}
        \label{fig:benchmark_scatter_containerlab_multi}
    \end{subfigure}
    \caption{Total tokens vs.\ run duration per individual run (Containerlab, $n=10$ per model).}
    \label{fig:benchmark_scatter_containerlab}
\end{figure}

Figure~\ref{fig:benchmark_scatter_containerlab} plots total token consumption against run duration
for every individual run.
In single-agent mode, \texttt{qwen3:30b} scatters to the upper-right (high tokens, mid-to-high
duration), while in multi-agent mode its cluster shifts to the lower-left (few tokens, lower
duration), reflecting the benefit of task decomposition.
\texttt{claude-opus-4-7} remains consistently in the low-duration region across both configurations.
\texttt{gpt-oss:120b} occupies the far-right in both plots, confirming its consistently high
duration regardless of configuration.
The horizontal dashed line for the Static Scenario reflects its fixed 29\,s execution with no
token consumption.
The spread within each model's cluster reveals run-to-run variability in how deeply the
\acrshort{AI} agents chose to enumerate hosts and services before concluding.


\begin{figure}[H]
    \begin{subfigure}{0.5\linewidth}
        \centering
        \includesvg[width=1\linewidth]{figures/benchmark_services_25_05_2026.svg}
        \caption{Single-agent.}
        \label{fig:benchmark_services_containerlab_single}
    \end{subfigure}%
    \begin{subfigure}{0.5\linewidth}
        \centering
        \includesvg[width=1\linewidth]{figures/benchmark_services_26_05_2026.svg}
        \caption{Multi-agent.}
        \label{fig:benchmark_services_containerlab_multi}
    \end{subfigure}
    \caption{Average hosts, services, and findings discovered per model (Containerlab, $n=10$ runs each).}
    \label{fig:benchmark_services_containerlab}
\end{figure}

Figure~\ref{fig:benchmark_services_containerlab} compares discovery breadth across models and
configurations.
\texttt{claude-opus-4-7} consistently discovers all 5 reachable hosts, 9 services, and the maximum
of 8 findings in both configurations, matching the deterministic baseline on findings.
\texttt{qwen3:30b} improves from 7 to 8 services in multi-agent mode, suggesting that sub-agent
specialisation compensates for weaker individual reasoning depth.
\texttt{gpt-oss:120b} similarly improves from 6 to 7 services and from 4 to 5 findings between
configurations.
The Static Scenario's 13 services exceed all \acrshort{AI} models because the script issues
exhaustive port sweeps unconditionally, while \acrshort{AI} agents terminate scanning early once
they judge the network to be sufficiently enumerated.


\subsection{Qualitativ Criteria}\label{subsec:qualitativ-criteria-containerlab}
